\section{Crossing scales and rigidity of local passage costs}
\label{met:section}

The metric argument is first carried out for the bilateral word.  We use
the joint local observations of Theorem~\ref{ker:local-kernels}, their
re-expression under changes of time unit, and the simultaneous annulus
estimate of Section~\ref{ann:section}.  All costs in this section are
numbers of original primal edges.  Put
\[
 p=\frac1{d_\gamma},\qquad \xi=\gamma p,\qquad
 Q=\frac2\gamma+\frac\gamma2>2.
\]
The reference metric is the intrinsic LQG metric $D=D_h$, with the
normalization fixed in the statement of the theorem.
Theorems~\ref{met:rigidity} and~\ref{met:normalizers} identify the local
passage law and its deterministic normalization.

The comparison strategy follows Gwynne and Miller's LQG metric uniqueness
proof: optimal constants, repeated annular opportunities and local
finite-energy field changes bring a geodesic near a shortcut and improve
the comparison
\cite[Section~1.5 and Sections~3--6]{GwynneMillerMetric}. The continuum
metric theory rests on LFPP tightness
\cite{DingDubedatDunlapFalconetTightness}, weak-metric estimates
\cite{DubedatEtAlWeakMetric}, and that existence and uniqueness theorem.
Here the objects being compared are port passages. Their local kernels,
selected changes of law and eventual endpoint completion require the
additional arguments below.  We continue to work in the original
oriented peanosphere coordinate, with its real and imaginary field
laws.  The canonical coordinate and loop identifications already
established do not change those laws.  This section identifies local
costs up to a deterministic subsequential scalar; Section~\ref{comp:section}
uses the intrinsic GH input to determine that scalar and reach every
vertex.

\subsection{Reference metric facts and annular cutoffs}

The reference metric is already defined on the continuum surface.  We
recall the estimates used to compare it with discrete passages; none
of these statements identifies a graph-distance limit.


Write $D_h(\cdot,\cdot;U)$ for the internal metric on $U$, and fix one
deterministic normalization of the $\gamma$-LQG metric.  Set
$p=1/d_\gamma$ and $\xi=\gamma p$. For a length metric $D$ and a
continuous function $f$, the Weyl weighted metric means
\[
 (e^{\xi f}\!\cdot D)(z,w)
 =\inf_P\int_0^{\len(P;D)}e^{\xi f(P(t))}\,\dd t,
\]
where the infimum is over $D$-rectifiable paths from $z$ to $w$,
parameterized by $D$-arclength \citep{GwynneMillerMetric}.
Write $h_r(0)$ for the circle average of $h$ on $\partial B_r(0)$.

\begin{theorem}[LQG metric inputs]\label{enc:lqg-metric}
The subcritical LQG metric is a field-measurable continuous length metric
inducing the surface topology.  Its internal metric on $U$ is determined
by $h|_U$.  It obeys
\begin{align}
 D_{h+f}&=e^{\xi f}\!\cdot D_h,
 &\mu_{h+f}&=e^{\gamma f}\mu_h,\label{enc:weyl}\\
 D_{h\circ\psi+Q\log|\psi'|}(z,w)
 &=D_h(\psi(z),\psi(w))\label{enc:covariance}
\end{align}
for continuous $f$ and conformal coordinate maps $\psi$.
The whole-plane GFF metric is proper and geodesic, and the geodesic
between each fixed pair of distinct points is almost surely unique.

For every compact planar set and every
$0<\beta<\xi(Q-2)$, there are almost surely finite $C$ and positive
$r_0$ such that, for $|z-w|<r_0$ in that set,
\begin{equation}\label{enc:internal-holder}
 D_h(z,w;B_{2|z-w|}(z))\le C|z-w|^\beta.
\end{equation}
The corresponding estimates before almost-sure truncation have
polynomially high probability uniformly over the reference scale, with
the normalization $r^{\xi Q}e^{\xi h_r(0)}$.
\end{theorem}

\begin{proof}
Existence, uniqueness, locality, and Weyl scaling are
\cite{GwynneMillerMetric}; conformal covariance is
\cite[Theorem 1.3]{GwynneMillerCovariance}.  Properness is \cite[Lemma~3.8]{DubedatEtAlWeakMetric}; a proper
length space is geodesic, as also noted at the start of Section~4
of that paper.
Fixed-pair uniqueness is \cite[Theorem 1.2]{MillerQianGeodesics}.
The precise internal estimate, including its probability and
normalization assertions, is
\cite[Lemma 3.20, equations (3.66)--(3.67)]{DubedatEtAlWeakMetric}.
Apply it along dyadic radii and Borel--Cantelli to obtain
\eqref{enc:internal-holder}.  The scale constant is
$r^{\xi Q}$ for the uniquely normalized LQG metric.  On a fixed compact
chart, local field absolute continuity and Weyl scaling transfer the
almost-sure assertion to our other real-field laws away from their
insertions.  Uniform probability bounds under a changed law require
the density control established when that law is used.
\end{proof}

\begin{lemma}[Deterministic geodesic anchors]\label{enc:geodesic-anchors}
Fix a deterministic countable dense subset $A$ of the plane.  Almost
surely every nonempty open subarc of every nonconstant whole-plane
$D_h$-geodesic contains a nondegenerate segment of the unique geodesic
between some two points of $A$.
\end{lemma}

\begin{proof}
Intersect fixed-pair uniqueness over $A\times A$ with the simultaneous
strong-confluence event of
\cite[Theorem 2 and Proposition 6]{BhatiaKavvadiasConfluence}.
In a given open geodesic subarc choose a strictly interior interval
$[s,t]$ and $0<\varepsilon<(t-s)/2$.  Proposition 6 supplies metric
neighborhoods of its two endpoints such that all geodesics between
those neighborhoods contain the original segment
$[s+\varepsilon,t-\varepsilon]$.  Since $A$ is also metric dense,
choose distinct anchors in these neighborhoods.  Their unique
geodesic contains the required segment.  The event is simultaneous
over all starting geodesics; only the anchor pair is selected afterward.
For later local applications, one first confines a sufficiently short
interior piece below the positive cost of exiting a surrounding buffer.
Locality and local absolute continuity then transfer this whole-plane
statement to that piece.
\end{proof}

The simultaneous annulus theorem can retain the following metric
regularity events together with a passage test.  Its field and spatial
alignment are unchanged.

\begin{corollary}[Additional local cutoffs]\label{ann:pairings}
The good events in Theorem~\ref{ann:aligned-annuli} may include
finitely many local internal \(D_h\) across and around regularity
tests, and bounds for the Gaussian pairings associated with any
fixed finite family of smooth Cameron--Martin recipes supported
strictly inside \(E\).  Their cutoffs can be chosen before applying
the theorem so that their combined marginal failure is arbitrarily
small, uniformly in the label.
\end{corollary}

\begin{proof}
By Theorem~\ref{enc:lqg-metric}, in each fixed buffered geometry
the internal metric induces the Euclidean topology, positive
separations between disjoint compact ports are positive, and the
required confined connection costs are finite.  Finite maxima,
positive minima, and moduli of continuity therefore admit
high-probability finite cutoffs and positive lower cutoffs.
Only internal metrics in the stated dependency buffer are used.
Under \eqref{ann:normalized-field}, the coordinate and constant
identities for this reference metric give
\[
 D_{h^{c,r,v}}(u,u')
       =v^{-p}D_h(c+ru,c+ru')
\]
with the corresponding internal domains.  These are the ordinary
unit-coordinate regularity events at every label.

For \(f\in C_c^\infty(E^\circ)\), the pairing
\((h,f)_\nabla\) is a centered Gaussian of variance
\(\|f\|_\nabla^2\) under the real reference law.  Constants pair
to zero; Dirichlet norms and pairings are invariant under affine
coordinate rescaling.  The same holds for the centered imaginary
field.  Gaussian tails and a finite union bound give any required
common probability cutoff.  On \(|(h,f)_\nabla|\le T\), the
Cameron--Martin factor
\(\exp((h,f)_\nabla-\|f\|_\nabla^2/2)\), and its reciprocal,
have fixed finite bounds.  These are field events and are adjoined
to the local output tests without changing their kernels.
\end{proof}


\subsection{Port comparisons and deterministic chaining}

We first specify the comparisons that will be proved.  A passage
certificate includes its ordered endpoint and intermediate ports, a
compact confinement, its upper cost, and the upper costs of its ordered
subpassages.  A certificate is a retained finite ordered path record, with arbitrary
positive losses in its port, confinement and cost requirements; it is
not merely the spatial range of a path.  Lemma~\ref{ker:readouts} and
Theorem~\ref{ker:local-kernels} retain these records jointly, including
their consistent refinements and restrictions.  Statements about
an open confinement are read using compact subsets of that confinement;
endpoint and intermediate ports can be made arbitrarily small.

\begin{definition}[Local passage comparisons]
\label{met:port-comparison}
Fix a joint limiting law of the field, its reference internal metrics,
and the compatible ordered passage records.  An \emph{upper comparison with constant
$C$} means the following simultaneous assertion.  For every
$D$-rectifiable path $P$ compactly contained in an open set $U$, every
finite ordered subdivision of $P$, and every positive endpoint, tube,
and cost tolerance, there is an attainable certificate which follows
the subdivided path through those tolerances, with each marked
subpassage costing at most $C$ times the corresponding $D$-length plus
the assigned cost tolerance.  An \emph{upper constant} is a deterministic
$C$ for which this holds almost surely in every ordinary chart.

A \emph{lower comparison with constant $c$} means that every finite-cost
certificate, with limiting endpoints $x,y$, has cost at least $cD(x,y)$.
It applies to every ordered subcertificate as well.  A \emph{lower
constant} is a deterministic $c\geq0$ for which these assertions hold
almost surely in every ordinary chart.

For $x\ne y$ in $U$, let $F_U(x,y)$ be the infimum of the upper costs of
certificates from $x$ to $y$ compactly confined in $U$, with the endpoint
ports shrinking to $x,y$.  Set $F_U(x,x)=0$.
\end{definition}

The upper comparison permits the endpoints of a discrete path to be
chosen within the prescribed ports.  It gives no bound on the cost of
reaching a prescribed lattice vertex.  This distinction persists until
Section~\ref{comp:section}.

Here is the deterministic use of a simultaneous annulus cover.  The
annuli have fixed shape parameters: a central hole, a surrounding
annular band, and a strictly larger buffer.  Their centers range over
fixed fine meshes at finitely many radii.  A positive number $w_i$ is
attached to each annulus; in applications it is $v_i^p$.  All shape
parameters are held fixed while a discrete approximation tends to its
limit.

\begin{lemma}[Chaining through an annulus cover]
\label{met:chain}
Suppose every point of a compact set $K$ belongs to a central hole of a
finite family of buffered annuli contained in an open set $U$.  Their
largest outer radius is at most $r_*$, their holes have a fixed positive
clearance from their bands, and their reference weights are $w_i>0$.
Assume that a reference path from a hole to its surrounding separator
has $D$-length at least $a w_i$, where $a>0$.
\begin{enumerate}
\item If the $i$th annulus contains an attainable surrounding primal
circuit of cost at most $Lw_i$, then every $D$-rectifiable path in $K$
can be followed, with endpoint errors tending to zero with $r_*$, at
cost at most
\[
 \frac{L}{a}\operatorname{len}_D(P)+L\max_i w_i.
\]
For any fixed finite ordered subdivision of the reference path, the
following walk can also retain its intermediate ports.  Each marked
subpassage has cost at most $(L/a)$ times the corresponding reference
length plus $2L\max_iw_i$.  All spatial errors are $O(r_*)$, and
arbitrary fixed open tube losses are allowed.
\item Suppose instead that every primal traversal from the hole side
to the outside of the band costs at least $\ell w_i$.  Suppose also
that a fixed reference circle lies strictly outside the tested band
and strictly inside its buffer, with internal $D$-diameter at most
$Aw_i$ in that buffer.  Every bounded-cost
certificate from $x$ to $y$ whose range lies in $K$ then satisfies
\[
 \operatorname{cost}\geq \frac{\ell}{A}D(x,y)-o_{r_*}(1).
\]
The error depends only on the reference metric modulus on the fixed
compact neighborhood of $K$.
\end{enumerate}
Both conclusions hold for limits of discrete paths when the annular
tests have strict margins and all selected ordered subrecords are
retained.
\end{lemma}

\begin{proof}
Fix the annulus cover before taking a discrete approximation.  Use the
strictly buffered upper tests to choose all its simple primal circuits
simultaneously on that approximation, with costs at most $Lw_i+o_n(1)$,
and denote their filled sides by
$\Omega_i$.  Transfer them to the reference sphere by $H_n$.  For large
$n$, each specified central hole lies strictly inside its $\Omega_i$,
and $\Omega_i$ lies in the corresponding outer disk.  The errors of
$H_n$ are smaller than the fixed clearances.  In the estimates below we
may first use any $a'<a$ to absorb those errors, and let $a'\uparrow a$
after taking the limit.

Starting at $P(0)$, select a hole containing that point and follow $P$
until its first encounter with the boundary of the selected filled
circuit.  At that encounter select another hole containing the point,
and repeat, stopping when the endpoint is reached.  If an index is
chosen several times, retain its occurrences separately.  Write
$\Omega_1,\ldots,\Omega_m$ for this ordered list and $[s_j,s_{j+1}]$
for the corresponding path intervals.  Each interval except possibly
the last has $D$-length at least $a'w_j$, and its image lies in
$\Omega_j$.  These intervals have disjoint interiors, so
\[
 \sum_{j<m}w_j\le (a')^{-1}\operatorname{len}_D(P).
\]
The finite cover has a positive minimum weight, so there are only
finitely many complete intervals, bounded uniformly in the approximation
by $1+\operatorname{len}_D(P)/(a'\min_iw_i)$.  Thus all summed
$o_n(1)$ cost errors vanish at this fixed cover.  The point $P(s_{j+1})$ belongs to
$\partial\Omega_j\cap\operatorname{int}(\Omega_{j+1})$.  Consequently
$\Omega_{j+1}$ cannot be contained in $\Omega_j$.

We reduce this list in its original order.  Keep an ordered list of
retained disks.  When a new disk contains the last retained disk,
delete that last disk, assigning its path interval to the new disk,
and repeat until no such containment remains.  The new disk cannot
be contained in the preceding retained disk: otherwise the last disk
just deleted would have been contained there as well, contrary to
the defining property of the retained list.  The new disk also overlaps the preceding disk in its interior, since
it contains the deleted disk whose interior overlapped that preceding
disk.
Thus consecutive retained disks have intersecting interiors and
neither contains the other.  Their Jordan boundaries therefore
intersect.  These are intersections of actual primal circuits, so
they contain common primal vertices.  No transverse-crossing claim
about limiting circuit traces is needed.

Choose a common vertex for each consecutive pair.  On each retained
circuit follow one arc between its two selected vertices; on the
first and last circuit choose the free endpoints arbitrarily.  If
there is only one retained circuit, both endpoints may be chosen
there.  The resulting primal walk uses at most the full cost of
each retained circuit, and hence at most
\[
 L\sum_{j=1}^m w_j
 \le \frac{L}{a'}\operatorname{len}_D(P)+L\max_iw_i.
\]
Each retained disk has spatial diameter $O(r_*)$, contains all path
intervals assigned to it, and the assignments remain consecutive in
path order.  Thus the walk's endpoints are within $O(r_*)$ of the
prescribed endpoints, and its circuits lie in any fixed open tube
around $P$ once $r_*$ is sufficiently small.

This ordered reduction also retains finitely many intermediate ports.
Assign a marked reference time to the retained disk containing its
path interval.  The entire corresponding circuit is within $O(r_*)$
of its marked reference point.  Choose each marked vertex to be the incoming intersection on its
circuit occurrence, or the free initial vertex on the first one.
Several marks assigned to one disk may use the same vertex; their
order is then preserved without an additional traversal.  For an individual subpassage,
all retained circuit intervals strictly between its two end marks
are charged to disjoint reference intervals within that subpassage.
At most the two circuits carrying the end marks are not so charged.
Its cost is therefore at most
\[
 \frac{L}{a'}\operatorname{len}_D(P|_{[t,t']})
       +2L\max_iw_i.
\]
The full-path bound above still has its single terminal term.  In
applications the number of prescribed marks is fixed and the covers
are refined with $\max_iw_i\to0$, so these boundary errors fit any
assigned positive subpassage tolerances.  All intersections and joins
were made on one discrete surface.  Pass to the retained finite
records only after this construction, then let $a'\uparrow a$ and
refine the covers.

For the lower bound, apply the same stopping construction to a discrete
path, stopping on the fixed circles beyond the tested bands.  Thus
every complete step contains a full band traversal.  Mark in every complete
step its traversal of the tested band.  The marks occur on disjoint
ordered path intervals.  A budget $T$ allows at most
$T/(\ell\min_i w_i)$ complete steps; thus they can be retained under
subextraction.  If $I$ is the set of selected complete steps, then
\[
 \operatorname{cost}\geq \ell\sum_{i\in I}w_i.
\]
Apply the same chronological deletion to the filled intermediate
reference circles, retaining a chain of intersecting boundaries.
Connect the successive circle-intersection points within their
respective buffers, using the internal-diameter bounds.  This gives
a reference route of length at most
$A\sum_{i\in I}w_i$.  The first and last reference connections have
Euclidean displacement at most a fixed multiple of $r_*$; their costs
are $o_{r_*}(1)$ by uniform continuity of $D$ on the compact
neighborhood.  Hence
\[
 D(x,y)\leq A\sum_{i\in I}w_i+o_{r_*}(1).
\]
The last, possibly incomplete, disk contributes only an endpoint error:
its reference circle is approached at the last complete step, and its
remaining Euclidean displacement is bounded by its diameter.

For discrete approximations, keep the finite cover, the strict band
margins, and all stopping marks fixed first.  The number of marks is
bounded by the displayed budget.  The subrecord property transfers the
traversal lower tests, and the attained-witness property transfers the
circuit upper tests.  Robust separation and intersection hold before
passing to a limit.  Finally refine the annulus covers.  This order of
limits does not require any estimate for costs from arbitrary lattice
endpoints.
\end{proof}

\begin{lemma}[Confined reference connections]
\label{met:rect-access}
Almost surely, on each compact subset of an ordinary chart there are
$\chi>0$, $H<\infty$, and $r_0>0$ such that
\begin{equation}
 D(x,y;B_{2|x-y|}(x))\leq H|x-y|^\chi,
 \qquad |x-y|\leq r_0.
 \label{met:internal-holder}
\end{equation}
For rectangles of bounded aspect ratio this gives a uniform internal
connection modulus up to all sides and corners.  The same bound, with a
constant depending only on the aspect ratio and the compact cutoff,
holds in every open neighborhood of the closed rectangle, uniformly as
that neighborhood decreases to the rectangle.
\end{lemma}

\begin{proof}
For a whole-plane field, Lemma~3.20 of \cite{DubedatEtAlWeakMetric}, equations
(3.66)--(3.67), gives the internal-ball bound with any
$0<\chi<\xi(Q-2)$, with polynomially high probability along dyadic
scales.  Borel--Cantelli gives \eqref{met:internal-holder}; local field
absolute continuity transfers its almost-sure version to the ordinary
charts used here.  Quantitative cutoffs are taken in the aligned
reference units before this transfer.

For completeness, let $R$ be a closed rectangle with bounded aspect
ratio, $z\in R$, and let $F_t(z)$ clamp the two coordinates of $z$ to
the rectangle whose sides are shifted inward by $t$.  Put
$t_n=(31/32)^nt$.  Then
\[
 |F_{t_n}(z)-F_{t_{n+1}}(z)|\leq \frac{\sqrt2}{32}t_n,
 \qquad \dist(F_{t_n}(z),\partial R)\geq t_n.
\]
The ball used in \eqref{met:internal-holder} therefore lies strictly
inside $R$.  Connecting successive clamped points and summing a
geometric series gives an interior path from $z$ to $F_t(z)$ of length
at most $H_1t^\chi$, including a continuous extension at $z$.  For
$x,y\in R$, use $t$ equal to the smaller of $|x-y|$ and one quarter of
the shorter side.  The segment between $F_t(x)$ and $F_t(y)$ lies in the
$t$-interior of $R$.  Divide it into steps of size at most $t/8$ and use
\eqref{met:internal-holder} on each step.  Bounded aspect ratio bounds
the number of steps when $|x-y|$ is comparable to the side length;
otherwise it is bounded absolutely.  Adding the two inward paths
proves the asserted modulus.  Every constructed path lies in the
closed rectangle, so enlarging the rectangle by an arbitrarily thin
open collar cannot increase its cost.
\end{proof}

\begin{lemma}[Vanishing joins and stability]
\label{met:vanishing-joins}
A finite upper comparison implies the following assertions,
simultaneously on each compact ordinary chart.
\begin{enumerate}
\item Every point has arbitrarily small surrounding primal circuits
whose limiting costs tend to zero, uniformly over the compact set.
\item The functions $F_U$ are symmetric, satisfy the triangle
inequality, and are locally continuous in their endpoints.  They obey
$F_U(x,y)\leq C D(x,y;U)$.
\item Nonconstant certificates can be concatenated at a common
limiting endpoint with arbitrary open and cost loss.  Comparisons with
fixed constants, strict shortcut witnesses, and saturation on compact
geodesic arcs pass to joint tangent extractions that retain the
reference internal metrics on their collars.
\end{enumerate}
Finite connected systems of reference ports and rectifiable paths can
in particular be followed by primal paths with all prescribed robust
intersections.  These assertions concern free endpoint ports.
\end{lemma}

\begin{proof}
Choose a fixed sufficiently fine angular subdivision of a small
Euclidean circle.  Lemma~\ref{met:rect-access} connects consecutive
vertices inside balls contained in its annular neighborhood.  The
result is a rectifiable reference closed path of winding one, with
length at most $H_2r^\chi$.  The same construction works uniformly for
centers on fine meshes and, by strict hole margins, surrounds every
point of the compact set.  Follow its successive pieces using the
upper comparison, with transverse overlapping ports.  The topology
comparison of Corollaries~\ref{top:transverse} and~\ref{top:circuits}
gives a surrounding primal walk and hence a simple
surrounding subcircuit, with cost at most
$CH_2r^\chi+o(1)$.  Take countably many gridded covers and then let
$r$ decrease.

To join two nonconstant certificates ending at $z$, choose a small
surrounding circuit which both certificates must meet.  Cut them at
their encounters with the circuit and join along that circuit.  The
added cost tends to zero.  This proves the triangle inequality; cases
with a constant subpassage use $F_U(z,z)=0$.  Symmetry comes from
reversing primal paths.  Following internal reference paths proves
$F_U\leq CD(\cdot,\cdot;U)$.  Consequently, for endpoints in a compact
subdomain,
\[
 |F_U(x,y)-F_U(x',y')|
 \leq C\bigl(D(x,x';U)+D(y,y';U)\bigr).
\]
Lemma~\ref{met:rect-access} gives the local continuity asserted here.

In a tangent extraction the reference field and its internal metric
are retained with their fixed joint law.  Their compact moduli are
therefore tight.  A strict low-cost witness uses finitely many compact
ports and retains its cost slack.  Conversely, a strictly cheaper
certificate in a proposed saturated limit has a finite attained
record in a strict collar.  In preceding laws, correct its endpoint
errors by the circuits just constructed.  Their costs tend uniformly
to zero on compact reference-modulus cutoffs, contradicting the
preceding saturation inequalities.  First fix the strict collar and
finite record, next take the extraction limit, and finally decrease
the endpoint tolerances.  This proves the stated stability without
requiring costs to prescribed lattice endpoints.
\end{proof}

\begin{lemma}[Marked strict improvement]
\label{met:improvement}
Consider a sequence of increasingly fine simultaneous annulus covers
with fixed reference regularity cutoffs $a,A>0$.  Include the tests of
Lemma~\ref{met:chain}, a lower bound $aw_i$ for reference advance
across a fixed thinner middle band, and lower bounds $aw_i$ for the
internal reference distance between opposite face ports of the finite
rectangle library constructed below.  The ports have strict disjoint
thickenings and the rectangles have strict surrounding collars.
Fix this finite library and then these cutoffs before applying the
simultaneous annulus theorem.  Corollary~\ref{ann:pairings} permits
their combined failure probability to be arbitrarily small; decreasing
$a$ and increasing $A$ incorporates the original chaining cutoffs.
\begin{enumerate}
\item Suppose $C$ is an upper constant.  If every marked reference
geodesic traversal in the good bands admits an alternative certificate
at ratio at most $C-\delta$, then $C-\varepsilon_+$ is an upper
constant for some $\varepsilon_+>0$ depending only on
$\delta,a,A$ and the fixed band geometry.
\item Suppose $c$ is a lower constant.  If every marked rectangular
subtraversal in the good bands costs at least $c_1$ times its internal
reference endpoint distance, where $c_1>c$, then
$c+\varepsilon_-$ is a lower constant for some
$\varepsilon_->0$ depending only on $c_1-c,a,A$ and the band geometry.
The second assertion permits $c=0$.
\end{enumerate}
It is enough that the indicated good annulus events have the high
probability and positive label density required by the simultaneous
annulus theorem.
\end{lemma}

\begin{proof}
Use the ordered stopping construction of Lemma~\ref{met:chain} and
mark a crossing of a thinner middle band in each complete step.  Its
reference advance is at least $a_1w_i$ for a fixed $a_1>0$.
The marked intervals are disjoint.  Connecting the outer reference
circles as in that lemma gives
\begin{equation}
 D(x,y)\leq A_1\sum_iw_i+o(1)
 \label{met:weight-cover}
\end{equation}
with a fixed $A_1<\infty$.

For the upper assertion, run the construction along a reference
geodesic.  Replace each marked interval by its improved certificate,
and retain the upper comparison $C$ on the intervening pieces.
At this fixed cover there are only finitely many gaps.  The complete
ordered upper comparison follows them with assigned cost tolerances
whose sum is arbitrarily small; Lemma~\ref{met:vanishing-joins}
supplies the joins with equally small total error.  These tolerances
are chosen after fixing the marked cover, so no error proportional
to its possibly growing number of marks is left when the cover is
refined.  The saving is at least
$\delta a_1\sum_iw_i$.  Use \eqref{met:weight-cover} and then refine
the covers.  For instance any
$0<\varepsilon_+<\delta a_1/A_1$ is admissible after the strict losses.
For a general rectifiable path in a prescribed open tube, its compact
range has positive reference distance from the complement of a
slightly smaller buffered tube.  Partition it into pieces shorter
than that distance.  Reference geodesics joining the consecutive
endpoints then stay in the tube and have total length no greater
than the original path.  Apply the preceding construction to those
pieces, keeping the specified ordered ports.  This proves the
complete upper-following assertion.

For the lower assertion, fix a normalized spatial length $h>0$ much
smaller than the thinner band's remaining crossing distance and its
clearance from the retained band collar.  Use squares of half-side
$h$ centered on a grid of mesh $h/4$.  The starting point of each
marked crossing lies within $h/8$ in each coordinate of one such
center.  Every square so selected stays in the retained band collar,
and the crossing must exit it.  Suppose its first exit is on the
right side $x=c_x+h$.  Retain the subpassage after its last crossing
of $x=c_x+h/2$ before that exit.  It stays in the exact rectangle
\[
 [c_x+h/2,c_x+h]\times[c_y-h,c_y+h]
\]
and joins its opposite vertical faces.  The other three exit sides
give the corresponding rectangles.  These form a finite library of
aspect-four rectangles; corner exits are allowed.  On discrete
approximations, round the stopping locations to nearby vertices and
use the strict face-port and collar margins, then retain their limits.
The ordered subpassages remain disjoint.  The stipulated finite-library
cutoffs give reference endpoint distance at least $a_2w_i$ in the
retained collar, with fixed $a_2>0$.  Apply the lower comparison $c$
on every gap and the improved comparison on each marked rectangle.
The reference triangle inequality for these finitely many ordered
endpoints yields
\[
 \operatorname{cost}\geq cD(x,y)+(c_1-c)a_2\sum_iw_i.
\]
If $c=0$, the strictly positive rectangle charge itself bounds the
number of complete steps at a fixed mesh.  If $c>0$, the same bound
holds and the lower comparison can also be subdivided intrinsically.
Thus the required marked subextractions exist in either case.
Equation~\eqref{met:weight-cover} gives any
$0<\varepsilon_-<(c_1-c)a_2/A_1$ after the strict losses.

Finally, the simultaneous annulus theorem makes these constructions
available about every path location with probability tending to one
along the specified covers.  Countable exhaustion of paths by ports,
then of chart compacts and tolerances, makes the improved comparison
an almost-sure statement.  No rate of discrete convergence is used:
the finite cover is always fixed before its discrete approximation.
\end{proof}

\subsection{Quantile scales and a first-exit lower bound}

For the normalization experiment represent the cone in its
circle-average embedding: its marked points are $0,\infty$ and
\[
 \sup\{r>0:h_r(0)+Q\log r=0\}=1.
\]
This is the convention of \cite[Definition~4.10]{DMSMating}.
Fix three strictly nested versions of one round annular circuit test,
with one fixed dependency buffer strictly inside $\mathbb D$ and
away from $0$.  At every microscopic time unit use this same embedded
cone input law, the same band geometries, and the same reference
mark convention.  In particular the field's additive height is fixed
by this convention; it is not reset within a selected local sample.
At time unit $4^j$, let $X_j$ be the least unscaled cost of the
middle test, using the contour correspondence and short-arc
interpolation to read winding.  On the event that the correspondence
does not yet resolve the strict buffers, use any positive finite
fallback value.  That event has probability tending to zero.  Fix a
quantile level $\tau<1$ sufficiently close to one, and define
\begin{equation}
 A_j=\inf\{t>0:\PP[X_j\leq t]\geq\tau\}.
 \label{met:quantile}
\end{equation}
Changing finitely many initial values makes every $A_j$ positive and
finite.  The topological comparison makes the circuit tests available
with probability tending to one; hence the fallback convention has no
effect on any conclusion below.  At atoms we use inequalities with
strict cost and probability slack.  In particular,
\[
 \PP[X_j<A_j]\leq\tau,\qquad
 \PP[X_j\leq A_j]\geq\tau.
\]
The narrower and wider bands allow the corresponding implications to
pass through the local certificate topology.

We shall use first exits from finitely many strict probability tests.
Here is their precise convention.  A bad traversal test consists of a
finite compact endpoint/confinement specification and a bounded cost
threshold.  Insert strict larger endpoint ports, a strict larger
confinement, and a strict larger cost threshold.  In the compactified
observation space, a continuous function with values in $[0,1]$ can be
chosen to equal one on the smaller bad test and zero outside the larger
bad test.  Work first on compact reference-data cutoffs and then
exhaust those cutoffs.  Tests of the expectation of this function are
continuous in the joint law.  A weak inequality is consequently
preserved under extraction, whereas a limiting strict inequality is
eventually preserved in the approximants.  The attained-witness and
subrecord assertions in Theorem~\ref{ker:local-kernels} are what make
the two buffered tests correspond to the required passage statements.
This convention will also be used for circuit tests and numerical
ratio tests.

\begin{lemma}[Initial growth]
\label{met:initial-growth}
The quantile scales satisfy
\[
 \limsup_{j\to\infty}\frac{\log_4 A_j}{j}\geq p.
\]
\end{lemma}

\begin{proof}
Fix $\varepsilon\in(0,p)$ and a sufficiently high-probability lower
traversal test in a buffered band.  Start at $j_0$ with cost unit one
and, at level $j\geq j_0$, use the denominator
\[
 t_{j,j_0}=4^{(p-\varepsilon)(j-j_0)}.
\]
At each fixed offset $j-j_0$, a bounded number of primal edges makes
vanishing spatial advance as $j_0\to\infty$.  Thus the traversal
lower test passes, with any fixed cost and probability slack, on
arbitrarily long initial runs.

Suppose first failures existed for arbitrarily large starts.  Their
distance from their starts tends to infinity.  Extract the joint local
observations at those failures, retaining every fixed preceding
offset, with denominator $t_{j,j_0}$.  At offset $i$ the passing lower
test gives a barrier proportional to
$4^{-(p-\varepsilon)i}$ in the current units.  The reference metric
unit there is $4^{-pi}$.  Their ratio is therefore $4^{\varepsilon i}$.
On labels $i\in[N,2N]$, the simultaneous annulus theorem and the lower
part of Lemma~\ref{met:chain} give a lower comparison at least
$c_0 4^{\varepsilon N}D$, with $c_0>0$ independent of $N$.
Letting $N$ tend to infinity says that every fixed separated traversal
has infinite cost in the extracted law.  It strictly passes the
original lower test.  The continuous buffered-test convention then
contradicts that these scales were failures.

Consequently a sufficiently late start has no later failure.  A
circuit in the normalization band contains a traversal of one of the
strict lower test bands.  The lower probability cutoff was chosen
stronger than needed for the quantile level in
\eqref{met:quantile}; hence, for all sufficiently large $j$,
\[
 A_j\geq c_\varepsilon 4^{(p-\varepsilon)j}
\]
with some $c_\varepsilon>0$.  This eventual lower bound will also ensure that a single edge has
vanishing cost in $A_j$ units.  Letting $\varepsilon$ decrease to zero
proves the asserted limsup.
\end{proof}

\begin{lemma}[A finite comparison from preceding scales]
\label{met:finite-upper}
Consider a joint extraction in units $A_j$ for which
\begin{equation}
 \limsup \frac{A_{j-i}}{A_j}\leq K4^{-pi}
 \quad\hbox{for every fixed }i\geq0
 \label{met:preceding-upper}
\end{equation}
with one finite $K$.  Its limiting passage law has a finite, positive
optimal upper constant $C$.  It has an optimal lower constant
$c\in[0,C]$.  Both constants are deterministic and attained in the
sense of Definition~\ref{met:port-comparison}.  Their inequalities
persist after every fixed change of volume unit and under further
joint tangent extractions retaining the reference internal metrics.
\end{lemma}

\begin{proof}
The quantile test at offset $i$, transported to its ordinary field
chart, supplies a high-probability circuit bound proportional to
$A_{j-i}/A_j\leq K4^{-pi}$.  Take the quantile sufficiently high that
the simultaneous annulus theorem applies after the fixed field
density changes and reference regularity cutoffs.  Its weights are
$w_i=4^{-pi}$.  The upper part of Lemma~\ref{met:chain} gives a finite
upper constant.  If the optimal upper constant were zero, follow a
reference circuit in the narrower normalization band at arbitrarily
small cost and stitch its finitely many transverse pieces.  The
reference circuit length is finite with arbitrarily high probability.
The middle normalization test would then cost less than, say, $1/2$
in $A_j$ units with probability greater than $\tau$, contradicting
\eqref{met:quantile}.  Thus $C>0$.

Define $C$ as the infimum of deterministic almost-sure upper constants
and $c$ as the supremum of deterministic almost-sure lower constants.
Use countably many rational ports, chart compacts and tolerances in
these definitions.  Intersect the probability-one events for constants
decreasing to $C$ and increasing to $c$; the compact certificate
extraction and arbitrary positive cost losses show that the limiting
constants themselves hold.  Symmetry and the existence of finite
nonconstant upper passages imply $c\leq C$.

For a volume factor $v$, the same discrete map with time unit $v4^j$
and denominator $A_jv^p$ has field input
$h-\gamma^{-1}\log v$ and costs $v^{-p}$ times the original costs.
This is the joint array re-expression of
Proposition~\ref{ker:field-transfer}.  It does not identify two differently
indexed kernels.  The ordinary field laws on a strict patch are
equivalent under the fixed shift, so the probability-one comparisons
retain the same constants in these units.  Further tangent
extractions retain their inequalities by
Lemma~\ref{met:vanishing-joins} and the ordered lower subrecords.
\end{proof}

\subsection{Optimal constants and a saturated tangent}

Fix a deterministic disk containing the compact annular geometries
and their collars used in this argument.  We now use a new ordinary
comparison experiment: the real field is a whole-plane GFF pinned by
its circle average on a deterministic circle outside that disk.  The
imaginary field keeps its independent whole-plane law and uniform
phase, and the local outputs use the same kernels.  Mutual local
field equivalence transfers the upper and lower comparisons on each
strict patch; a countable atlas makes them simultaneous.  Their
optimal constants are therefore the same $C,c$.  This choice precedes
all witness and hit selections.  No realized field is repinned.
The later neighborhood $U$ and, after discarding finitely many coarse
windows, all modification supports will lie inside the fixed disk.
Consequently every modifying function vanishes on the chosen pin.

We also fix the distinction between the law, its indexed coordinates,
and a further tangent.  A unit test at label $v$ evaluates its indexed
kernel under this one ordinary normalized input law, the same for
every $v$.  Thus $h_v$ in
\eqref{ker:units} has that fixed law; $h=h_v+\gamma^{-1}\log v$
expresses the old quantum units.  We do not extract the diverging
fields $h-\gamma^{-1}\log v$ for one fixed unnormalized $h$.
Fixed finer offsets are retained jointly from this normalized input.
The almost-sure comparison constants transfer to each such test by
the indexed identity and fixed-label field equivalence.  The uniform witness probabilities below are proved directly under
this common ordinary input law, using the annulus theorem and
optimality of the constants.  A varying-label tangent is a new joint
limit of those unconditioned indexed laws, formed with the normalized
input held fixed as in Proposition~\ref{ker:normalized-diagonal}.
Auxiliary shortcut draws will be sampled only after that tangent has
been fixed; they do not replace its original passage system.

We will prove that the constants in Lemma~\ref{met:finite-upper}
coincide.  For the moment suppose $c<C$.  A nearly saturated witness at
volume label $v$ is a reference geodesic traversal of a fixed closed
annular band for which every competing certificate in a fixed strict
open enlargement costs at least $(C-\delta)$ times its length.  A
shortcut witness is a certificate in a closed rectangle $R$, whose
endpoints lie on opposite faces, with cost less than $c_1$ times their
reference internal distance in a fixed open collar $\Omega$ of $R$.
The rectangle comes from a fixed gridded library in the annular band.
Endpoint separation, collar clearance and strict cost losses are
included in the observations.

We use the absence of these witnesses only with a reserved ratio
slack.  For a fixed geodesic traversal, failure of the lower bound on
all competing certificates means that one finite attained record has
cost below that bound; rationally enlarged ports and confinement and
a rational cost threshold still witness the strict inequality.
Conversely, absence of all such shortcut records means that every
retained rectangle subcertificate obeys the asserted lower bound.
The reference internal metric and the ordered record sets are both
part of the local observation.  On fixed length and modulus cutoffs,
the reference geodesic family is compact in uniform parametrization;
its geodesic identities are checked at rational times.  The endpoint
continuity in Lemma~\ref{met:vanishing-joins} and countable port
refinement make the corresponding witness events measurable.  In
applying their negations below, first reserve an intermediate ratio,
then fix the port and collar losses, and only afterward take an
extraction or refine an annulus cover.  Thus the negations supply the
finite alternatives and universal lower subrecords used in
Lemma~\ref{met:improvement}.

\begin{lemma}[Many labels carry extremal witnesses]
\label{met:extremal-labels}
There are fixed band geometries and $\eta>0$ such that the following
holds.  For every $\delta>0$ and all sufficiently large $N$, more than
three quarters of $i\in[N,2N]$ have a nearly saturated witness at
$v=4^{-i}$ with probability at least $\eta$.  For each fixed
$c_1\in(c,C)$, the analogous statement holds for shortcut witnesses.
The lower probability and spatial separation for nearly saturated
witnesses can be chosen independently of $\delta$.
\end{lemma}

\begin{proof}
Choose the band geometry and high-probability reference regularity
cutoffs first.  Choose $\eta$ below the failure threshold in the
simultaneous annulus theorem for a label set of density one quarter.
If the first assertion failed along arbitrarily large windows, at
least one quarter of their labels would have no nearly saturated
witness with probability greater than $1-\eta$.  On these labels every
marked reference geodesic traversal has an alternative at ratio below
$C-\delta$, with the retained strict open loss.  The first part of
Lemma~\ref{met:improvement} decreases the upper constant by a positive
amount, contradicting its optimality.  The probability threshold used
in this argument depends on the annular geometry and its reference
regularity, not on $\delta$; only the required fineness and cost losses
depend on $\delta$.

If the shortcut assertion failed, a quarter of the labels would, with
the requisite high probability, have no low-ratio rectangle witness.
Every marked rectangular subpassage would then cost at least
$c_1D(\cdot,\cdot;\Omega)$.  The second part of
Lemma~\ref{met:improvement} strictly increases the lower constant.
This is a contradiction also when $c=0$.  Start with an intermediate
ratio between $c$ and $c_1$ if a strict collar or cost loss must be
absorbed.  The finite rectangle construction in that lemma supplies
the fixed positive spatial separation.  Reference diameter and
modulus cutoffs, taken with probability loss less than $\eta/2$, make
all retained lengths finite and separated from zero in the normalized
units.
\end{proof}

\begin{lemma}[Saturation with deterministic anchors]
\label{met:dense-anchors}
If $c<C$, a further tangent law has upper constant exactly $C$,
admits strict shortcuts with ratio bounded away from $C$, and has the
following property with positive probability.  There are fixed
distinct deterministic anchors, rational $0<t_1<t_2<1$, and a fixed
bounded open set $U$ such that the unique reference geodesic $P$
between those anchors, parametrized on $[0,1]$ proportionally to
$D$-length, satisfies
\begin{equation}
 F_U(P(t_1),P(t_2))
   =C D(P(t_1),P(t_2)),\qquad P([t_1,t_2])\Subset U.
 \label{met:saturated}
\end{equation}
The event can be restricted to deterministic positive lower and
finite upper bounds on the total and interval lengths.
\end{lemma}

\begin{proof}
Let $\delta_k\downarrow0$.  The two sets of labels in
Lemma~\ref{met:extremal-labels} intersect.  Choose indices $i_k\to\infty$
in these intersections and put $v_k=4^{-i_k}$.  This choice is deterministic,
since it uses probabilities under the common normalized input law.
Extract the unconditioned joint laws, including every fixed finer
offset and the reference internal metrics, by the restricted diagonal
construction of Proposition~\ref{ker:normalized-diagonal}.  The discrete indices
are chosen so that the effective time unit $n_kv_k$, and every
retained fixed multiple of it, tend to infinity.  Only after this
joint extraction do we retain the positive-probability witness events.
On a fixed reference-modulus and
length cutoff, the nearly saturated geodesics have equicontinuous
arclength parametrizations, lie in the fixed compact band, and have
separated endpoints.  Extract their paths and the internal metric on
their fixed open collar.  A strictly cheaper limiting certificate
would be attained with strict ports in the preceding laws, and its
endpoint errors could be corrected by
Lemma~\ref{met:vanishing-joins}.  This contradicts the near-saturation
inequality.  The limiting arc is therefore saturated.  The same
extraction keeps the fixed-gap shortcut witness by its strict cost
and spatial margins.  Its upper optimal constant is $C$, since the
inherited upper bound holds and saturation forbids a smaller one.
Its lower optimal constant is strictly below $C$; applying
Lemma~\ref{met:extremal-labels} in this law supplies a fixed positive
density of still finer labels with shortcut ratio at most
$C-\delta_0$, for one $\delta_0>0$.

Take a short strictly interior subarc of the saturated arc.  It is an
ambient reference geodesic: when an internal collar was used to take
the extraction, shorten the arc until its length is below the
positive cost of leaving a still larger collar.  We are in the exterior-pinned whole-plane experiment fixed above.
Proposition~6 of \cite{BhatiaKavvadiasConfluence} applies simultaneously
to every reference geodesic.  It gives endpoint neighborhoods such
that geodesics joining points in those neighborhoods contain a
nondegenerate interior segment of our arc.  Shrink these endpoint
neighborhoods so that short internal connections to the arc, followed
by the arc, cost less than the positive cost of leaving the larger
collar.  Every joining geodesic then stays in that collar.  Thus the
anchor selection uses only the protected internal metric, as required
for the local field-law transfer.  Choose the endpoints from a fixed
countable Euclidean dense set.  They are metric dense as
well.  Fixed-pair uniqueness holds simultaneously for the countably
many pairs by \cite[Theorem~1.2]{MillerQianGeodesics}.  Thus the saturated arc contains
a nondegenerate segment of one of these unique anchor geodesics.

Saturation passes to its subarcs: otherwise splice a cheaper
subpassage with the upper comparison on the two remaining pieces and
contradict saturation of the larger arc.  Choose rational fractional
times inside this segment and a rational open neighborhood with
compact clearance.  Countably many anchor pairs, rational intervals,
neighborhoods and length cutoffs cover the positive-probability
saturation event.  One fixed choice has positive probability, giving
\eqref{met:saturated}.
\end{proof}

\subsection{Amplifying and preserving a shortcut}

We work in the tangent of Lemma~\ref{met:dense-anchors}.  Its finer
favorable labels have a shortcut ratio at most $C-\delta_0$, where
$\delta_0>0$ is fixed.  Here $v=4^{-i}$ is a favorable shortcut label,
$\widehat v=Mv$ is the displaced annular label in
Theorem~\ref{ann:aligned-annuli}, and $u=v^p$.  By the
reference cutoffs in Lemma~\ref{met:extremal-labels}, its denominator
$d=D_{h}(z_0,z_1;\Omega)$ can be required to satisfy
\begin{equation}
 a'u\leq d\leq A'u,
 \qquad \operatorname{cost}(\mathcal S)\leq(C-\delta_0)d
 \label{met:shortcut}
\end{equation}
for constants $0<a'<A'<\infty$.  The certificate $\mathcal S$ lies in a
closed rectangle $R$ and has endpoints on opposite faces.  A
preservation patch $V$ contains $R$ with strict margin and satisfies
$\overline V\subset\Omega$.  We also require
\begin{equation}
 D_h(z_0,z_1;V)\leq A'u
 \label{met:preservation-access}
\end{equation}
uniformly as the open rectangular collar $V$ decreases to $R$.
Lemma~\ref{met:rect-access} bounds the internal diameter of $R$
in every such collar.  Its normalized cutoff has arbitrarily high
probability, uniformly in the label.  Increase $A'$ and restrict the
unit shortcut test by this cutoff, losing less than half its fixed
positive probability.  Thus \eqref{met:preservation-access} and the
small terminal-cap moduli are fixed before $M$ is chosen.  After a
real shift the same requirements are read for $D_{h+b}$.

\begin{lemma}[Many compartments]
\label{met:amplification}
Fix a desired shell failure probability $\varepsilon>0$ and a
sufficiently large fixed $M$.  Put $w=(Mv)^p$.  There are finitely
many disjoint buffered compartments inside a fixed annular shell and
a finite deterministic set of real shifts $\mathcal B$ with the
following property, uniformly over the favorable labels.  With
probability at least $1-\varepsilon$, one compartment and one
$b\in\mathcal B$ have an auxiliary passage draw satisfying
\eqref{met:shortcut} for $D_{h+b}$, in actual cost units $u$, with the
specified preservation patch, the bound
\eqref{met:preservation-access} for $D_{h+b}$, and the reference
regularity cutoffs.
The event is local in the union of the compartment buffers.  The
necessary traversal order can be required to be determined in their
fixed larger imaginary-field buffers.
\end{lemma}

\begin{proof}
First restrict the positive-probability unit shortcut test so that its
traversal order is determined inside a fixed finite imaginary buffer.
The finite-order localization and whole-plane flow-tree determination
in Lemma~\ref{ker:finite-order} allow the omitted probability to be
less than half the original success probability.  Fix this buffer
before packing its copies.  The resulting unit test has probability
at least some $\eta_0>0$.

Inside the chosen shell, pack $k$ smaller copies with mutually
disjoint full buffers.  Their relative shape is fixed; their spatial
scale decreases when $k$ increases.  Decompose both fields into fresh
Dirichlet parts and exterior harmonic parts on those buffers.  After
subtracting the coarse real value and the imaginary phase, the
harmonic oscillation on a strict inner buffer has a scale-uniform
Gaussian bound.  Choose a bound $H$ so that, with probability at least
$1-\varepsilon/4$, at least $k/2$ compartments have the required
harmonic bound.  Indeed the expected fraction failing any one bound
is uniformly small, so Markov's inequality gives this conclusion;
independence of the harmonic parts is not used.

Here is the exact coordinate and volume compensation.  If the parent
annulus has center $c$, radius $r$, and label $Mv$, its normalized
real input is
\[
 g(x)=h(c+rx)+Q\log r-\gamma^{-1}\log(Mv).
\]
For a compartment $x=d+az$, with $a>0$, its input at label $v$
after the preservation shift $b$ is
\begin{equation}\label{met:compartment-input}
 g(d+az)+Q\log a+\gamma^{-1}\log M+b.
\end{equation}
Write $H_c$ for its exterior-measurable coarse value in $g(d+a\cdot)$.
Round $-H_c-Q\log a-\gamma^{-1}\log M$ to the fixed mesh for $b$.
The residual is bounded uniformly; the absolute label $v$ has
canceled.  This chooses the field of a new auxiliary draw.
By \eqref{ker:units}, $O_v=v^{-p}O$ on the same incidence array,
and the reference metric has the identical scalar re-expression.
Thus a unit-index shortcut has actual cost $u$ times its unit cost,
with $u=v^p$, and its preservation kernel is exactly the kernel for
$h+b$.  No change-of-field rule for a fixed passage cost is used.

On a compartment whose harmonic profile is bounded by $H$, bounded
Cameron--Martin extensions and the bounded residual shift give a
uniform conditional lower probability $p_0>0$ of success, as in
Lemma~\ref{ann:conditional-transfer}.  To see the
direction of this estimate, restrict the reference harmonic profile
to a fixed bounded set still carrying positive success probability.
Between that restriction and any allowed target profile the Gaussian
likelihood ratio has a uniformly bounded second moment.  If $L$ is
the likelihood ratio and $E$ the retained unit event,
\[
 \PP_{\rm ref}(E)^2
 \leq \PP_{\rm target}(E)\,
       \EE_{\rm target}[L^2].
\]
This bounds the target probability below uniformly.  The same
calculation applies to the imaginary harmonic part modulo its phase.
Conditional passage kernels do not change the likelihood ratio.
The constants depend on $\eta_0,H$, the fixed buffer shapes and the
residual bound, and not on the absolute size of the compensating
shift or on $k$.

Now choose $k$ so large that $(1-p_0)^{k/2}<\varepsilon/4$.
Conditional on the exterior fields, the fresh fields and the passage
draws in these disjoint buffers have product laws.  Thus all good
compartments fail with probability at most this amount.  With this
finite geometry fixed, truncate the coarse compensation values to a
bounded range, losing probability less than $\varepsilon/4$.
The possible mesh values now form the finite set $\mathcal B$.
For each $b$, sample its full joint kernel coordinate, retaining the
passage records, reference metrics, all required volume offsets, and
feasible-order outputs under the shifted local input.  Given the
fields, different prospective $b$ coordinates may be sampled
independently; their within-coordinate consistency is preserved.
The disjoint-compartment arrays have product conditional laws.
All these draws precede selection of a successful coordinate.
The finite-order data and their directed local traversals are
retained at the same time.  Additional reference
regularity cutoffs on these finitely many fixed geometries can be
chosen to lose the remaining $\varepsilon/4$.

This selection order is essential: the profile bound fixes $p_0$
before the number of compartments is chosen; the finite shift menu is
truncated afterward.  Uniformity comes from the bounded residual and
the exact array re-expression, not from continuity of passage costs
in a real-field shift.
\end{proof}

The next lemma is deterministic once the reference metric and shortcut
have been given.  For each proposed smooth function $f$, suppose a
passage system has upper comparison $C$ with $D_{h+f}$ and contains
the retained shortcut on $V$, where $f=b$.  These are its only
required properties.  Write $F^f$ for its infimum certificate cost
with compact confinement in the shell.  The lemma constructs $f$
and obtains a strict gain in that system.  Lemma~\ref{met:transfer-dominance} subsequently
constructs such an output system with the required dominated law.

\begin{lemma}[A finite family of two-gate modifications]
\label{met:gate}
Fix the cutoff data of Lemma~\ref{met:amplification}, including its
finite exact shift menu.  Work in shell coordinates in which all
modifications are supported in
$B_2(0)\setminus\overline{B_{1/8}(0)}$.  Let the reference anchor
geodesic enter $B_{1/4}(0)$ and let both anchors lie outside $B_2(0)$.
Write $l_-,l_+$ for the least reference costs from the anchors to
$\partial B_2(0)$.  On the reference shell cutoff suppose
\begin{equation}
 D(z_-,z_+)-l_--l_+\geq aw.
 \label{met:landing-gap}
\end{equation}
Choose $M$ so that $A'u<aw/8$.  There is a finite deterministic family
of smooth functions, compactly supported in the permitted shell, with
uniformly bounded amplitudes and Dirichlet norms, such that whenever a
shifted shortcut from Lemma~\ref{met:amplification} is available, one
of these functions has the following effect.

It equals the shortcut's exact shift $b$ on its preservation patch
$V$.  For the modified reference metric
$D^f=D_{h+f}=e^{\xi f}\cdot D_h$, every geodesic between the anchors
contains a segment compactly within the shell for which that passage
system admits a competing certificate confined in the shell, of cost
strictly less than $C$ times its reference length.  The saving is at
least $\delta_0a'u/2$.  The certificate uses the preserved shortcut on
$V$ and the upper comparison for its two approaches.  All inequalities admit strict spatial and cost margins.
The probability lost in imposing the required reference moduli and
finite-domain cutoff bounds can be made arbitrarily small before
using the simultaneous annulus theorem.
\end{lemma}

\begin{proof}
The proof concerns the reference metric and a retained passage
certificate.  A Weyl identity is used only for $D^f$.
Choose
\[
 0<e<\min\left\{\frac{A'u}{100},
                    \frac{\delta_0a'u}{16(C+1)}\right\},
 \qquad B_0=2A'u.
\]
All individual approximate-path and gluing errors below are allocated
smaller fixed fractions of $e$.

\emph{Landing and terminal caps.}
Let $x_-,x_+\in\partial B_2(0)$ minimize the two exterior landing
costs.  One can equivalently use approximate minimizers within the
reserved error.  For landing sets within $\tau$ of those minima,
the reference triangle inequality and
\eqref{met:landing-gap} give mutual distance at least $aw-2\tau$.
The reference upper modulus therefore separates these sets by a
uniform positive Euclidean amount on the cutoff event.
Choose two disjoint outer caps around $x_-,x_+$, with all original
connection errors less than $e$.  Around the shortcut endpoints
$z_0,z_1$ choose disjoint terminal caps whose internal $D_{h+b}$
connection errors to the respective endpoints are less than $e$ in
$\Omega$.  Choose outer stops just inside $\partial B_2(0)$ and
terminal stops just outside $V$.  Decrease the margin of $V$ below the
terminal-cap scale.  Lemma~\ref{met:rect-access} provides the uniform
connection estimates needed for this operation, including when an
endpoint lies at a corner.

\emph{The allowed regions; see Figure~\ref{met:gate-figure}.}
Join the outer stops to their corresponding terminal stops by two
disjoint simple polygonal arcs outside $V$, avoiding the central
omitted disk and the opposite caps.  To construct them, first join
the outer boundary component to the rectangular boundary component.
Cutting along this arc leaves the planar domain connected, so the
second prescribed joining arc can be drawn there.  Fix their short
end approaches before thickening the arcs.

Use two nested thickenings to obtain a smaller allowed set $E_0$ and
a larger set
\[
 E_1=O\cup T_0\cup S\cup T_1.
\]
Here $O$ contains the exterior and a thin inner collar of
$\partial B_2(0)$; the chamber $S$ contains $V$ with
$\overline S\subset\Omega$; the tubes $T_0,T_1$ are disjoint;
$O\cap S=\varnothing$; and the only overlaps are
\[
 O\cap T_j\subset\mathcal A_j,
 \qquad S\cap T_j\subset\mathcal B_j,
 \qquad j=0,1,
\]
where $\mathcal A_j$ and $\mathcal B_j$ are the outer and terminal
caps.  The relevant overlaps have closures in slightly larger caps
with the same error estimates.  The smaller set has positive
clearance from the complement of the larger set along every interior
boundary.  An outer wall layer lies between them, compactly in the
permitted support, so exiting $E_1$ from the intended route requires
crossing a fixed positive-width wall.

Inside each tube take a closed core $K_j$, a finite union of
overlapping rectangles of positive width, connecting its two stops.
It is separated from $O$, $V$, and the wall layer.  Its old internal
diameter is bounded by $Kw$ on a fixed-domain cutoff.  A transverse
marker slab spans the whole tube away from the caps.  Its two
longitudinal faces have a fixed positive separation.  Every complete
tube transit crosses this slab.

\emph{A smooth function with the required plateaux.}
After fixing this geometry, let $m_*w>0$ be a lower cutoff for a
crossing of any wall or marker slab.  Choose $H\geq\max\{0,b\}$ and
$J\geq\max\{0,-b\}$ such that
\[
 e^{\xi H}m_*w>2B_0,
 \qquad e^{-\xi J}Kw<e/3.
\]
Next choose $\rho>0$ so small that a point within $\rho$ of a core
can access that core at old cost less than $e e^{-\xi H}/3$ in a
slightly larger local collar.  Keep these core collars disjoint from
$O,V$, and the wall.  The choice order is $H$, then $J$, then
$\rho$; the lower marker cutoff uses its fixed longitudinal faces
and does not depend on $\rho$.

Choose a smooth baseline $g$ which is nonnegative on $O$, at least
$b$ on $S$, exactly $b$ on $V$ and the smaller terminal caps, zero
on the smaller outer caps and near both support boundaries, and
exactly $H$ on the walls and markers.  It can be required to lie
between $\min\{0,b\}$ and $H$.  These closed plateau requirements
have positive clearance, so smooth cutoff functions construct $g$.
Let $\chi_j$ be smooth, equal to one on $K_j$, supported in its
$\rho$-collar, with the two supports disjoint.  Set
\begin{equation}
 f=(1-\chi_0-\chi_1)g-J(\chi_0+\chi_1).
 \label{met:gate-function}
\end{equation}
Then $f=b$ on $V$, $f=0$ outside the permitted shell, and $f=-J$
on each core.  It is nonnegative on $O$, at least $b$ on $S$ outside
the tubes, and equals $H$ on the wall.  On a marker it equals $H$
except within the $\rho$-collar of the core.  Outer cap hops have
multiplier at most one, and terminal cap hops have multiplier at
most $e^{\xi b}$.

\begin{figure}[htbp]
\centering
\begingroup
\definecolor{gatecore}{RGB}{30,106,79}
\definecolor{gatetube}{RGB}{236,246,240}
\definecolor{gatechamber}{RGB}{235,242,252}
\definecolor{gatehigh}{RGB}{250,226,211}
\definecolor{gatepath}{RGB}{29,75,132}
\begin{tikzpicture}[x=1cm,y=1cm,font=\small,
  cheap/.style={draw=gatetube,line width=13pt,line cap=round,line join=round},
  wall/.style={draw=gatehigh,line width=24pt,line cap=round,line join=round}]
  \node[anchor=west,font=\small\bfseries] at (0,4.90)
    {(a) Spatial placement inside the permitted shell};
  \fill[black!3] (2.60,2.35) circle (2.12);
  \draw[black!60,densely dashed] (2.60,2.35) circle (2.12);
  \draw[gatehigh,line width=8pt] (2.60,2.35) circle (1.86);
  \draw[wall] (1.45,3.41) .. controls (1.70,3.94) and (3.37,3.97) .. (3.65,2.96);
  \draw[wall] (1.45,1.29) .. controls (1.70,.76) and (3.37,.73) .. (3.65,1.74);
  \filldraw[fill=gatehigh,draw=none,rounded corners=7pt]
    (3.13,1.33) rectangle (4.18,3.37);
  \draw[cheap] (.94,3.57) .. controls (1.70,3.94) and (3.37,3.97) .. (3.65,2.96);
  \draw[cheap] (.94,1.13) .. controls (1.70,.76) and (3.37,.73) .. (3.65,1.74);
  \filldraw[fill=gatechamber,draw=gatepath!70,rounded corners=3pt]
    (3.27,1.48) rectangle (4.04,3.22);
  \filldraw[fill=white,draw=gatepath!70]
    (3.42,1.94) rectangle (3.90,2.76);
  \node at (3.66,2.35) {$V$};
  \node at (3.66,3.04) {$S$};
  \filldraw[fill=white,draw=black!65,densely dashed]
    (2.45,2.35) circle (.58);
  \node[align=center,font=\scriptsize] at (2.45,2.35)
    {omitted\\disk};
  \node[anchor=east] at (.76,3.71) {$\mathcal A_0$};
  \node[anchor=east] at (.76,.99) {$\mathcal A_1$};
  \node at (2.18,4.04) {$T_0$};
  \node at (2.18,.68) {$T_1$};
  \node[anchor=west] at (4.68,2.58) {$\partial B_2$};
  \node[anchor=west] at (4.68,2.15) {$O$ outside};
  \node[anchor=north west,align=left,text width=7.4cm,inner sep=0pt]
    at (6.70,4.25) {The two tubes join distinct outer caps to the
      chamber $S$ and avoid the central omitted disk.  The orange
      collar represents the high wall around the interior route.
      \par\medskip
      The wall and the field modification have positive clearance
      from both support boundaries.  The outer region $O$ also
      includes a thin inner collar of $\partial B_2$.
      \par\medskip
      This panel shows placement only; the adjacency and marker
      constraints are detailed below.};
\end{tikzpicture}

\par\medskip
\begin{tikzpicture}[x=1cm,y=1cm,font=\small,
  region/.style={draw=black!65,line width=.65pt,rounded corners=2pt,
                 minimum height=2cm,inner sep=0pt},
  caplink/.style={draw=black!75,line width=.8pt}]
  \node[anchor=west,font=\small\bfseries] at (0,3.18)
    {(b) Allowed-region adjacency};

  \node[region,fill=black!4,minimum width=1.6cm] (outerleft) at (.85,1) {};
  \node[region,fill=gatetube,minimum width=2.6cm] (tubezero) at (3.85,1) {};
  \node[region,fill=gatechamber,minimum width=3.1cm] (chamber) at (7.5,1) {};
  \node[region,fill=gatetube,minimum width=2.6cm] (tubeone) at (11.15,1) {};
  \node[region,fill=black!4,minimum width=1.6cm] (outerright) at (14.15,1) {};

  \draw[densely dashed,black!55,line width=.55pt]
    (outerleft.north) -- (.85,2.55) -- (14.15,2.55) -- (outerright.north);
  \node[fill=white,inner xsep=5pt,inner ysep=2pt] at (7.5,2.55)
    {same exterior region $O$, drawn twice};

  \draw[caplink] (outerleft.east) --
    node[above=4pt] {$\mathcal A_0$} (tubezero.west);
  \draw[caplink] (tubezero.east) --
    node[above=4pt] {$\mathcal B_0$} (chamber.west);
  \draw[caplink] (chamber.east) --
    node[above=4pt] {$\mathcal B_1$} (tubeone.west);
  \draw[caplink] (tubeone.east) --
    node[above=4pt] {$\mathcal A_1$} (outerright.west);

  \node at (.85,1.28) {$O$};
  \node at (.85,.72) {$f\geq0$};
  \node at (14.15,1.28) {$O$};
  \node at (14.15,.72) {$f\geq0$};
  \node at (3.85,1.60) {$T_0$};
  \node at (11.15,1.60) {$T_1$};
  \node at (7.5,1.68) {$S$};

  \filldraw[fill=gatecore!12,draw=gatecore,line width=.7pt,rounded corners=1pt]
    (2.70,.40) rectangle (5.00,.95);
  \node at (3.85,.675) {$K_0,\quad f=-J$};
  \filldraw[fill=gatecore!12,draw=gatecore,line width=.7pt,rounded corners=1pt]
    (10.00,.40) rectangle (12.30,.95);
  \node at (11.15,.675) {$K_1,\quad f=-J$};
  \filldraw[fill=white,draw=gatepath!75,line width=.65pt,rounded corners=1pt]
    (6.45,.30) rectangle (8.55,1.27);
  \node[align=center] at (7.5,.79) {$V$\\[-1pt]$f=b$};

  \node[anchor=north] at (7.5,-.22)
    {$\mathcal A_j$: outer caps \qquad $\mathcal B_j$: terminal caps};
\end{tikzpicture}

\par\medskip

\begin{tikzpicture}[x=1cm,y=1cm,font=\small,
  transit/.style={draw=gatepath,line width=1.05pt,-{Stealth[length=2mm]}}]
  \node[anchor=west,font=\small\bfseries] at (0,3.75)
    {(c) A transverse marker forces access to the core};

  \fill[gatetube] (.35,.25) rectangle (9.15,2.95);
  \fill[gatehigh] (4.10,.25) rectangle (5.80,2.95);
  \node[anchor=west] at (.48,3.20) {$T_j$};
  \draw[decorate,decoration={brace,amplitude=3pt},black!65]
    (4.10,3.06) -- (5.80,3.06);
  \node[anchor=south] at (4.95,3.15) {marker slab};

  \filldraw[fill=white,draw=gatecore!65,densely dashed,line width=.65pt,
            rounded corners=3pt]
    (.82,1.05) rectangle (8.68,2.10);
  \filldraw[fill=gatecore!14,draw=gatecore,line width=.8pt,
            rounded corners=1pt]
    (1.00,1.38) rectangle (8.50,1.76);
  \node at (2.40,1.57) {$K_j,\quad f=-J$};
  \node at (4.95,2.53) {$f=H$};
  \node at (4.95,.63) {$f=H$};

  \draw[black!55,line width=.55pt] (4.10,.25) -- (4.10,2.95);
  \draw[black!55,line width=.55pt] (5.80,.25) -- (5.80,2.95);
  \draw[black!65,line width=.7pt,rounded corners=2pt]
    (.35,.25) rectangle (9.15,2.95);

  \draw[transit] (.65,2.43)
    .. controls (2.35,2.56) and (3.32,1.98) .. (4.30,1.95)
    -- (5.70,1.95)
    .. controls (6.80,1.95) and (7.48,2.46) .. (8.85,2.43);
  \fill[gatepath] (4.55,1.95) circle (1.7pt);
  \node[anchor=south east,inner sep=2pt,text=gatepath] at (4.50,2.02) {$p_j$};
  \draw[gatecore,line width=.65pt,-{Stealth[length=1.4mm]}]
    (4.55,1.91) -- (4.55,1.61);
  \node[anchor=north west,inner sep=0pt] at (.84,.91)
    {core's $\rho$-collar};
  \draw[black!60,line width=.45pt] (2.47,.89) -- (2.88,1.05);

  \node[anchor=north west,align=left,text width=5.0cm,inner sep=0pt]
    at (9.72,2.96) {A complete transit with $D^f$-length at most $B_0$
      must enter the $\rho$-collar, since
      \[
        e^{\xi H}m_*w>2B_0.
      \]
      From such a point $p_j$, the cheap core and terminal hop give
      \[
        D^f(p_j,z_j;\text{local collar})<e.
      \]};
\end{tikzpicture}
\endgroup

\caption{The gate construction in Lemma~\ref{met:gate}.
(a) Schematic placement of the chamber and two tubes away from the
central omitted disk.  The orange collar indicates the high wall;
its exact shape and width are not prescribed by the drawing.
(b) Adjacency of the allowed regions; the two $O$ boxes denote the same
exterior region.  The only overlaps are in the outer caps $\mathcal A_j$
and terminal caps $\mathcal B_j$.  The preservation patch $V\subset S$
has $f=b$, and the tube cores $K_j$ have $f=-J$.
(c) On a transverse marker, $f=H$ outside the core's $\rho$-collar.
The crossing-cost cutoff forces a transit of $D^f$-length at most $B_0$
into that collar; the cheap core and terminal hop then give access to
$z_j$.  The high wall confines the route to $E_1$.}
\label{met:gate-figure}
\end{figure}

\emph{Reference forcing.}
Follow the original exterior minimizing paths, the cheap cores and
their cap hops, and an internal connection across $V$.  This gives a
modified competitor of length at most
\begin{equation}
 l_-+l_++A'u+2e<l_-+l_++B_0.
 \label{met:gate-budget}
\end{equation}
For any modified anchor geodesic, the portion between its first and
last visits to $\partial B_2(0)$ has length at most $B_0$, since the
two exterior portions cost at least $l_-,l_+$.  This local portion
cannot cross the high wall, and hence stays in $E_1$.  Every complete
tube transit must approach its core within $\rho$ while crossing the
marker; otherwise it pays more than $B_0$ at height $H$.
At such a marker point $p_j$, use the access to the core, its cheap
internal connection, and the terminal hop.  These give
\begin{equation}
 D^f(p_j,z_j;\hbox{local collar})<e.
 \label{met:gate-access}
\end{equation}

The first outer cap visited must be $\mathcal A_0$.  If it were
$\mathcal A_1$, the prefix lies in $O$, where its old length is no
larger than its modified length.  Joining its last point to $x_+$
inside that cap and following the old exterior path to $z_+$ would
give an old anchor connection of length at most
$l_-+l_++B_0+e$, contrary to \eqref{met:landing-gap}.
The last outer cap is similarly $\mathcal A_1$.

Take the first visit to $\mathcal A_1$ after the first visit to
$\mathcal A_0$, and then the last visit to $\mathcal A_0$ before it.
The intervening open path avoids both outer caps.  The allowed-set
overlaps imply that it lies wholly in $O$ or wholly in
$T_0\cup S\cup T_1$.  The first possibility again gives an old
competitor, now with at most two cap errors, of length
$l_-+l_++B_0+2e<l_-+l_++aw$.  Thus there is an inner transition
from $T_0$ through $S$ to $T_1$.

On that transition, take the first terminal-cap-$1$ visit and the
last terminal-cap-$0$ visit before it.  Between them the path lies
in $S$ outside both tubes.  Its shift is at least $b$, and the cap
endpoints are each within internal $D_{h+b}$ distance $e$ of their
respective $z_j$.  The central segment therefore has modified length
at least $d-2e$.  Select marker points $p_0,p_1$ on the complete tube
transits before and after this segment.  As they belong to a
geodesic,
\[
 D^f(p_0,p_1)\geq d-2e.
\]
This last/first selection incurs at most one error per cap regardless
of repeated excursions.

\emph{Passage gain and finiteness.}
Use \eqref{met:gate-access}, the output upper comparison, and the
retained shortcut certificate.  Concatenating with its strict
margins gives
\[
 F^f(p_0,p_1)
 \leq (C-\delta_0)d+2Ce
 \leq C D^f(p_0,p_1)-\bigl(\delta_0d-4Ce\bigr).
\]
The final parenthesis is at least $\delta_0a'u/2$, allowing the
reserved open losses.

Here is the order that makes the recipe family deterministic.  Fixing
$M$ and the amplification geometry fixes $e/w$ and $B_0/w$.  The
initial reference-modulus cutoffs give deterministic cap radii and
positive endpoint separations.  All endpoint and rectangle parameters
therefore lie in a deterministic compact family with fixed aspect and
boundary-clearance bounds.  Each valid polygonal construction persists
on a neighborhood of its parameters.  Choose a finite subcover of these parameter families;
equivalently use sufficiently fine deterministic endpoint and
rectangle bins.  First add uniform reference-modulus cutoffs for this
finite library, then its wall and core cutoffs, and finally choose
$H,J,\rho$ uniformly over it.  The resulting finite list of smooth
functions has bounded amplitude and bounded Dirichlet norm.
These later wall and core cutoffs change the plateau amplitudes; they
do not change the earlier spatial parameter family, gap or budget.
Translation and Euclidean dilation preserve the Dirichlet energy in
two dimensions.  All cutoff losses can be made arbitrarily small
without changing the earlier gap $aw$ or the budget $B_0$: the new
wall and core constants are absorbed by $H,J$.  This proves the
asserted finite, strict construction.
\end{proof}

\subsection{The weighted change-of-measure contradiction}

Fix the anchor pair, interval and neighborhood in
\eqref{met:saturated}.  Restrict the saturation event to total
reference length in a fixed compact subinterval of $(0,\infty)$ and
to a fixed positive clearance of the saturated interval in $U$.
All modifications below are centered near its middle third.  For a
favorable label $v$, use the shell event of
Lemmas~\ref{met:amplification} and \ref{met:gate}, aligned to $Mv$.
Its weight is $w=(Mv)^p$.  The original geodesic's passage through the
central hole supplies \eqref{met:landing-gap}.  Require also the
reference outer-circle diameter bound $Aw$ and the local lower
crossing cutoffs used in the gate construction.

For the finite list of smooth functions, add a cutoff on every
Gaussian pairing appearing in its Cameron--Martin likelihood.  The
functions have uniformly bounded Dirichlet norms after spatial
rescaling, so these finitely many pairing cutoffs have arbitrarily
high marginal probability, uniformly in the label.  The simultaneous
annulus theorem applies to their intersection with the other shell
events.

\begin{lemma}[Transfer before selection]
\label{met:transfer-dominance}
For each fixed ball, label and gate recipe, let $\mu$ be any
subprobability measure obtained by restricting the source joint law
to saturation, a successful hit, that recipe, and its likelihood
cutoff.  Its image under the one-ball modification can be completed
to a measure $\nu$ satisfying
\begin{equation}
 \nu\leq K\PP,
 \label{met:domination}
\end{equation}
where $\PP$ is one common unmodified joint law of the tangent field,
reference metric and passage observations.  The finite $K$ is
independent of the ball, label, window length and selected hit.
\end{lemma}

\begin{proof}
Fix the recipe before imposing a hit or saturation condition.
Condition on both fields.  On its preservation patch retain the
auxiliary draw for the exact field $h+b$, and outside a strictly
larger modification buffer retain the original exterior observations.
Their joint law is the product-kernel marginal in
Lemma~\ref{ker:selected-transfer}.  Under $h'=h+f$, these are
exactly the two required output kernels, since $f=b$ on the first
patch and $f=0$ on the second.  Complete all other output observations
with their correct conditional law.  The imaginary field is unchanged.

The only density change is now that of the real field.  The
Cameron--Martin density, evaluated on the source field, is
\[
 L_f(h)=\exp\left((h,f)_\nabla+\frac12\|f\|_\nabla^2\right).
\]
On the imposed likelihood cutoff it is at most $K$.  Therefore the
unselected image restricted to that cutoff is bounded by $K\PP$.
Restricting further to any event determined by the old observations,
the auxiliary shortcut, the hit or the chosen recipe gives a
submeasure and preserves this domination.  Thus no conditional
independence after selecting a hit is required.  The maximum cutoff
over the fixed finite recipe list is uniform under spatial dilation
and translation, giving the claimed single $K$.
\end{proof}

\begin{lemma}[Strict loss approaches equality]
\label{met:near-equality}
There are deterministic $\epsilon_N\downarrow0$ and source events
whose probabilities tend to one such that every selected
modification in the $N$th annular window satisfies, in its output,
\begin{equation}
 0<C-
 \frac{F_U(P'(t_1),P'(t_2))}
      {D_{h'}(P'(t_1),P'(t_2))}
 \leq\epsilon_N.
 \label{met:near-equality-event}
\end{equation}
Here $P'$ is the output reference geodesic between the same anchors.
The denominator remains bounded above and away from zero.
\end{lemma}

\begin{proof}
All supports have radius at most a deterministic number tending to
zero, their centers range over a fixed compact set, and their
amplitudes are bounded by a fixed constant.  The old $D$-diameters of
these supports tend uniformly to zero.  On either an old or a
modified geodesic, the part between the first and last support visits
can be replaced by a reference connection between two points of that
support, at cost bounded by a fixed amplitude factor times its old
diameter.  Outside the support the two length rules agree.
Consequently all modified anchor paths have old length tending to the
old anchor distance.  Global comparison by the bounded multiplier
keeps them in a fixed proper metric ball.  Their arclength
parametrizations are equicontinuous; every subsequential limit is an
old geodesic.  Uniqueness identifies it with $P$.  The collapsed
support-spanning portions have vanishing arclength, so fractional
times converge as well.  This reasoning is uniform over the finite
recipe geometries and all their centers in the compact set: a
failure of uniformity would give a contradicting sequence of
shrinking supports.

In particular, the gate segment lies inside the same fixed interval
$[t_1,t_2]$ of $P'$ for all sufficiently fine windows, with compact
clearance in $U$.  Apply the output upper comparison to the two
remaining pieces of this interval and use the strict gain from
Lemma~\ref{met:gate}.  This proves the strict left inequality in
\eqref{met:near-equality-event}.

For the other inequality, take a proposed output competitor in $U$
with bounded cost.  Cut it at its first and last encounters with a
slightly larger ball than the observation buffer.  Its exterior
subcertificates have the original costs and can be retained in the
old law.  Replace the whole intervening portion by one old short
connection, and correct its two moved endpoints by old short
connections.  Lemma~\ref{met:vanishing-joins} bounds all added costs
by a deterministic function of the old compact reference modulus,
the largest support radius, and the endpoint displacements, tending
to zero.  Repeated visits cost no additional error because only the
first and last encounters are used.  If the competitor never enters
the ball, it transfers directly.  Strict larger buffers ensure
compact confinement of the retained exterior records even when an
individual competing path has only its own positive clearance in
$U$.

An output competitor whose cost was below the old saturated value
by a fixed positive amount would therefore contradict
\eqref{met:saturated}.  This proves the required lower bound on its
cost up to an error tending uniformly to zero on old reference-data
cutoffs.  The reference endpoint distances converge, and the fixed
interval length has a positive lower bound.  The resulting ratio
error tends to zero in probability.  Choose a deterministic sequence
$\epsilon_N\downarrow0$ by successive probability cutoffs; retain
the source event on which all errors are at most $\epsilon_N$.
Its probability tends to one, proving the assertion.
\end{proof}

\begin{lemma}[Weighted image-count contradiction]
\label{met:weighted-count}
The saturated tangent of Lemma~\ref{met:dense-anchors} cannot admit
the fixed-gap shortcuts described there.
\end{lemma}

\begin{proof}
Let $\mathcal X$ denote the measurable space of the complete joint
observations, and let $\mathcal I_N$ be the finite set of centers, scales, volume
labels and recipes in the simultaneous annulus construction, and
write $w_i=(Mv_i)^p$.  These weights are deterministic.  Under the
one common unmodified output law $\PP$, define
\[
 R=\frac{F_U(P(t_1),P(t_2))}{D_h(P(t_1),P(t_2))},
 \qquad E_N=\{0<C-R\leq\epsilon_N\}.
\]
Set $R=+\infty$ unless $P([t_1,t_2])\Subset U$ and the denominator
is positive.  Thus $R$ is a single random variable, independent of
the choice of ball or recipe.

For an index $i$, let $G_i$ be the following event of the output
alone.  The output anchor geodesic satisfies the fixed total-length
cutoff and makes a complete transit of the recipe's marker slab
inside its enlarged shell, with subpath length at least $a_3w_i$.
Also require the alignment test, expressed using the output field
minus that recipe's fixed function $f_i$.  Choose $a_3>0$ below the
source marker-crossing cutoff times the minimum bounded Weyl factor.
Every modified good hit has these properties.  All marker faces,
shells and functions are deterministic for index $i$; the reference
geodesic is measurable and almost surely unique.  Hitting times,
with the retained strict spatial margins, therefore make $G_i$
measurable from the output field and geodesic.  It refers to no
discarded source path or auxiliary draw.

On the saturation event, intersected with the high-probability cover and
error events, every point of the middle third of the interval has
at least $cN$ good covering holes.  The time spent by a reference
geodesic in any one ball is at most $Aw_i$: its first and last
outer-circle visits can be connected within the surrounding buffer at this
cost.  Integrating the covering multiplicity along an interval of
length bounded below gives
\begin{equation}
 \sum_{i\in\mathcal I_N}w_i\mu_i(\mathcal X)\geq c_1N
 \label{met:source-weight}
\end{equation}
for $c_1>0$, where $\mu_i$ is the source submeasure for the
corresponding good hit and a selected valid recipe.  Choose one
valid recipe measurably from the fixed finite list for each hit.
The source probability of the retained saturation event stays
bounded below, since its original probability was positive.

Push each submeasure separately through its modification, obtaining
$\nu_i$.  By Lemma~\ref{met:transfer-dominance}, $\nu_i\leq K\PP$
for the same $\PP$ and $K$.  By Lemma~\ref{met:gate} and
Lemma~\ref{met:near-equality}, it is supported on the output event
$G_i\cap E_N$ defined above.

The output events satisfy
\begin{equation}
 \sum_{i\in\mathcal I_N}w_i\1_{G_i}
 \leq c_2N
 \label{met:output-weight}
\end{equation}
on the fixed output total-length cutoff.  At a fixed spatial radius,
choose the first qualifying marker transit specified by each $G_i$.
Its length is at least $a_3w_i$ and its whole parameter interval
stays in that enlarged shell.  The mesh balls have bounded overlap,
so integration along the output geodesic bounds the sum of these
subpath lengths by a fixed multiple of its total length.  There are $O(N)$ spatial radii.  The
alignment coarse value is measured outside the modification support;
alternatively, its change is bounded by the fixed amplitude.
Consequently only boundedly many volume labels can be aligned with
one spatial ball.  The finite number of recipes contributes another
fixed factor.  This proves \eqref{met:output-weight}.

Combining the two weight bounds and the common-law domination gives
\[
 c_1N
 \leq\sum_iw_i\nu_i(\mathcal X)
 \leq K\int\1_{E_N}\sum_iw_i\1_{G_i}\,\dd\PP
 \leq Kc_2N\PP(E_N).
\]
For a fixed random variable $R$, the probability
$\PP(0<C-R\leq\epsilon_N)$ tends to zero.  This is a
contradiction.  All modifications were applied separately; neither
independence of selected hits nor a simultaneous modification of
overlapping balls was used.
\end{proof}

\begin{theorem}[Rigidity of a finite port comparison]
\label{met:rigidity}
Every limiting law satisfying \eqref{met:preceding-upper} has
coincident optimal constants $c=C=\lambda$, where
$0<\lambda<\infty$ is deterministic.  In every ordinary open chart,
\[
 F_U(x,y)=\lambda D_h(x,y;U).
\]
The lower comparison applies to all certificates; the upper
comparison has the complete ordered-port meaning of
Definition~\ref{met:port-comparison}.
\end{theorem}

\begin{proof}
If $c<C$, Lemma~\ref{met:dense-anchors} constructs a saturated
tangent with fixed-gap shortcuts.  Lemma~\ref{met:weighted-count}
excludes it, so $c=C$.  Positivity and finiteness follow from
Lemma~\ref{met:finite-upper}.

The upper inequality for the displayed internal distance follows
by following internal reference paths.  For the reverse inequality,
start with a bounded-cost certificate compactly confined in $U$ and
its actual discrete approximating paths.  Before taking limits, add
ordered marks whenever their accumulated normalized cost first crosses
a multiple of $2^{-k}$, retaining the previously specified marks.
The single-edge cost tends to zero, since the initial-growth argument
gives $A_j\to\infty$.  For each fixed $k$ the number of marks is
bounded by the cost budget.  Lemma~\ref{ker:readouts} lets us retain
all these finite refinements consistently along a diagonal.

Apply the lower comparison to the subrecord between any two retained
marks.  Their $D$-distance is at most $1/c$ times their intervening
cost.  This is the modulus that was unavailable before $c>0$.
The marked points therefore extend to a continuous path parametrized
by accumulated cost; flat cost intervals have identical spatial
endpoints.  Its image stays in the retained compact confinement.
For every finite subdivision, the sum of its reference endpoint
distances is at most the certificate cost divided by $c$.  Taking
the supremum over subdivisions bounds its $D$-length by that value.
Thus the certificate costs at least $cD(x,y;U)$, and infimizing gives
the reverse inequality.  This construction uses consistent finite
refinements and the positive lower bound; no continuous path is
inferred from a spatial range alone.  Exhaust compact confinements and retain
the arbitrary positive open losses to obtain the assertion for
every open $U$.
\end{proof}

\subsection{Propagation of the deterministic normalizers}

\begin{lemma}[Quantile control in a rigid law]
\label{met:rigid-quantiles}
There are deterministic $0<k_-<k_+<\infty$, depending only on the
fixed quantile and buffered test geometries, with the following
property.  If a joint extraction in $A_j$ units has comparison
$\lambda D$, then every fixed offset $t\in\Z$ has
\begin{equation}
 k_-\lambda4^{pt}
 \leq \liminf\frac{A_{j+t}}{A_j}
 \leq \limsup\frac{A_{j+t}}{A_j}
 \leq k_+\lambda4^{pt}.
 \label{met:rigid-ratios}
\end{equation}
In particular $\lambda\in[k_+^{-1},k_-^{-1}]$.  The statement
includes forward offsets whose ratios were first extracted in
$[0,\infty]$.
\end{lemma}

\begin{proof}
Fix $t\in\mathbb Z$ and put $v=4^t$.  In the joint extraction retain
the raw observations at effective time unit $v4^j=4^{j+t}$ and cost
denominator $A_jv^p$; denote their local kernel array by $K_t$.
Retain the shifted-base-coordinate description and the prescribed
circle-average cone description jointly, as spatial relabellings
of the same time-labelled raw records in the common aggregate
reference extraction, before changing any field law.
This is the effective-time array supplied by
Theorem~\ref{ker:local-kernels}(1),(4).  Evaluating $K_t$ with the
fixed embedded cone input and band convention used in
\eqref{met:quantile} therefore gives the limiting buffered
observations of $X_{j+t}/(A_jv^p)$.  The time-labelled raw records
are those of that same word model at time $4^{j+t}$, and their spatial
relabelling uses this fixed input.  No identification of $K_t$ with
the kernel at time $4^j$ is asserted.

The indexed unit identity of Proposition~\ref{ker:field-transfer}
instead identifies $K_t$ at the shifted input
$h_v=h-\gamma^{-1}\log v$: its passage costs and reference metric
are both $v^{-p}$ times those at the original input $h$.
The comparison constant in this shifted-input experiment is thus
$\lambda$.  On the fixed dependency buffer, the shifted cone
restriction and the unshifted circle-average cone restriction are
mutually absolutely continuous.  The same kernel $K_t$ is used under
these two input laws.  Hence their probability-one local comparison
with constant $\lambda$ transfers to the fixed unshifted cone input.
This step uses only an almost-sure transfer for each fixed $t$;
it does not require a density bound uniform in $t$.  Internal metrics
and all circuit and traversal tests here are confined to the retained
buffer.  Countably many fixed offsets may be retained together.

Now take the reference cutoffs under that one fixed cone input law.
A circuit in the middle band contains a reference advance across
a strict smaller band.  Its positive lower cutoff can have probability
sufficient to contradict the upper quantile if the ratio in
\eqref{met:rigid-ratios} were too small.  The all-certificate lower
comparison gives the left bound.  Conversely a rectifiable reference
winding path in the narrower band has a finite-length cutoff with
probability greater than $\tau$, with the needed probability slack.
Following and stitching it gives a middle-band circuit and hence the
right bound.  Both cutoffs depend only on this fixed input law and
the test geometries, so the constants $k_-,k_+$ are independent of
$t$ and of the extracted law.  Strict buffered tests and the two
quantile inequalities after \eqref{met:quantile} justify these
transfers, including possible quantile atoms.  Setting $t=0$ gives
the stated compact interval for $\lambda$.  The same bounds exclude
zero or infinite forward ratios before renormalizing at that offset.
\end{proof}

\begin{lemma}[Comparable preceding power ratios]
\label{met:majority}
Suppose a joint extraction at scales $j\to\infty$ satisfies the
following tests at every fixed preceding offset: adjacent
$A$-ratios lie in a fixed compact subinterval of $(0,\infty)$,
there is the normalized upper circuit quantile, and a fixed
sufficiently high-probability normalized lower traversal test
holds.  Then
\[
 r_i:=4^{pi}\lim\frac{A_{j-i}}{A_j},\qquad i\geq0,
\]
can be extracted with $r_i\in(0,\infty)$, and one constant $L_0$
satisfies $L_0^{-2}\leq r_i\leq L_0^2$ for all $i$.
\end{lemma}

\begin{proof}
The adjacent bounds give positive finite limits for every fixed
offset after a diagonal extraction.  Fix $\theta<1/4$.  Choose a
large $L_0$, uniformly in the starting offset $i$.
If, for arbitrarily large $N$, more than $\theta N$ values of
$k\in[N,2N]$ satisfy $r_{i+k}>L_0r_i$, the preceding lower tests
at these deterministic labels give barriers of size at least a
fixed constant times $L_04^{-pk}$ in the start-$i$ units.
The simultaneous annulus theorem and
Lemma~\ref{met:chain} force a lower comparison proportional to
$L_0D$ at the starting scale.  For sufficiently large $L_0$ this
contradicts its upper circuit quantile, because a fixed reference
band crossing has positive distance with the required probability.
All constants are uniform in $i$: at the starting scale use its own
$A_{j-i}$ and evaluate its effective-time kernels with the same fixed
cone input and band convention as in \eqref{met:quantile}.  The
ordinary local laws used by the simultaneous annulus theorem are
likewise fixed at every label.  Thus the reference probability cutoffs
in this contradiction are chosen once; no uniform absolute-continuity
bound for an unbounded family of additive shifts is being used.

Likewise, if more than $\theta N$ labels have
$r_{i+k}<r_i/L_0$, their upper circuit quantiles and upper chaining
give a following constant proportional to $L_0^{-1}$.  A reference
traversal with a fixed finite-length cutoff then contradicts the
starting normalized lower traversal test.  Increase $L_0$ once to
handle both contradictions.  Consequently, for every fixed $i$
and all sufficiently large $N$, more than
$(1-2\theta)N$ indices $\ell\in i+[N,2N]$ satisfy
\begin{equation}
 r_i/L_0\leq r_\ell\leq L_0r_i.
 \label{met:majority-band}
\end{equation}

For two fixed starts $i,i'$, the shifted intervals differ in at
most $2|i-i'|$ indices.  Since $1-2\theta>1/2$, the two sets of
indices in \eqref{met:majority-band} intersect for sufficiently
large $N$.  At an intersection index,
$L_0^{-2}\leq r_i/r_{i'}\leq L_0^2$.  Take $i'=0$, for which
$r_0=1$.  The constant is independent of $i$, as required.
\end{proof}

\begin{theorem}[All-scale port normalization]
\label{met:normalizers}
There is a deterministic positive function $B(n)$, tending to
infinity, and constants $s>0$, $K<\infty$ such that
\begin{equation}
 \frac{B(n')}{B(n)}\leq K\left(\frac{n'}n\right)^s,
 \qquad 1\leq n'\leq n.
 \label{met:polynomial-ratio}
\end{equation}
One can take $s=p/2$ and $B(n)=A_{\lfloor\log_4n\rfloor}$,
after changing finitely many initial values.

Along every deterministic sequence of diverging time units, joint
crossing-data extractions in $B(n)$ units have local comparison
$\lambda D$ with a deterministic $\lambda$ in a fixed compact
subinterval of $(0,\infty)$, after further extraction.  The
comparison includes all finite-cost lower certificates, complete
ordered-port upper following, the simultaneous small surrounding
circuits of Lemma~\ref{met:vanishing-joins}, and finite connected
systems of paths with prescribed transverse intersections.
The same comparison holds on every ordinary compact patch of a
conditioned finite-sphere extraction using these deterministic
units and the common local kernels.
\end{theorem}

\begin{proof}
We propagate a finite list of tests.  Choose a broad adjacent-ratio
interval $(L^{-1},L)$ which is strictly passed by the ratios of
Lemma~\ref{met:rigid-quantiles}.  Put $s=p/2$ and choose an integer
$m$ so large that those rigid ratios strictly pass
\[
 \frac{A_{j-m}}{A_j}<4^{-sm}.
\]
Finally choose a sufficiently small fixed normalized lower
traversal threshold so that every rigid law with
$\lambda\in[k_+^{-1},k_-^{-1}]$ strictly passes its
high-probability lower test.  Use the continuous buffered-test
convention above, with enough probability slack that limits of
passing tests still suffice for the simultaneous annulus theorem.
The upper circuit quantile is already supplied at every index by
\eqref{met:quantile}.  The remaining finite test family consists of
the adjacent-ratio bounds, the displayed $m$-step bound, and the
buffered lower traversal test.  We use two different implications:
limits of passing tests retain the closed inequalities needed for
Lemma~\ref{met:majority}, whereas every rigid law passes the tests
with strict margins and therefore forces nearby approximants to
pass.  The first-failure argument below uses these implications in
that order.

There are arbitrarily long success runs at arbitrarily large
indices.  Indeed, by Lemma~\ref{met:initial-growth}, for every
$\varepsilon>0$ the sequence
$4^{-(p-\varepsilon)j}A_j$ has arbitrarily late record maxima.
Take record indices with $\varepsilon\downarrow0$.  For every
fixed preceding offset $i$ their ratios satisfy
\[
 \frac{A_{j-i}}{A_j}\leq4^{-(p-\varepsilon)i}.
\]
Every joint limit therefore satisfies
\eqref{met:preceding-upper}.  Theorem~\ref{met:rigidity} and
Lemma~\ref{met:rigid-quantiles} imply that all the chosen tests
strictly pass at every fixed backward and forward offset.
If some bounded window of tests failed along these record indices,
another joint extraction would contradict that strict passage.
A diagonal choice yields success runs whose lengths tend to
infinity.

Suppose there were infinitely many later failures.  Starting from
the success runs just obtained, extract at the level immediately
before the first later failure.  Every fixed preceding offset
passes the tests.  Lemma~\ref{met:majority} supplies the uniform
exact-power bound \eqref{met:preceding-upper}.  The extracted law
is rigid, and Lemma~\ref{met:rigid-quantiles} includes the first
forward offset.  Thus all the tests strictly pass at the alleged
failure as well.  The continuity and strict-loss convention for
the tests gives a contradiction.  Hence all sufficiently large
indices pass.

In particular, for those indices,
$A_{j-m}/A_j\leq4^{-sm}$ and the adjacent ratios are bounded.
Write $j-j'=km+r$, with $0\leq r<m$.  Iterating gives
\[
 \frac{A_{j'}}{A_j}\leq L^r4^{-smk}
 \leq L^m4^{sm}4^{-s(j-j')}.
\]
Changing the constant handles the finite initial indices.  Passing
from $j,j'$ to $\lfloor\log_4n\rfloor,
\lfloor\log_4n'\rfloor$ changes this bound by at most another
fixed factor, proving \eqref{met:polynomial-ratio}.
Taking $n'=1$ proves $B(n)\to\infty$.

For an arbitrary sequence of powers of four, all fixed preceding
tests now pass.  Extract their ratios and apply
Lemma~\ref{met:majority} again.  Theorem~\ref{met:rigidity}
identifies the local port law, and
Lemma~\ref{met:rigid-quantiles} puts its deterministic factor in
the fixed positive compact interval.  For a general time unit
$n$, extract also $n/4^{\lfloor\log_4n\rfloor}\in[1,4]$.
The bounded time-unit re-expression of the contour, field and
certificate arrays absorbs this converging factor into $\lambda$;
its possible $p$th powers remain in a fixed compact interval.
No convergence rate is needed for this extraction.

The simultaneous circuit and finite-system assertions follow from
Lemma~\ref{met:vanishing-joins} and the robust topological
intersection statements.  Finally, the finite-volume part of
Theorem~\ref{ker:local-kernels} uses the same kernels and the same
deterministic cost units.  Lemma~\ref{ker:sphere-density}
dominates the actual normalized sphere field, with its actual
additive height, by the ordinary field law on each such patch.
Together with Proposition~\ref{ker:field-transfer}, it transfers
the probability-one comparisons there.  A countable chart exhaustion makes them
simultaneous off the exploration root.

For the sphere with $n$ edges the duration-one letter clock has unit
$2n$, so this first gives the comparison in $B(2n)$ units.
The propagated adjacent-ratio bounds imply that $B(2n)/B(n)$ lies
in a fixed positive compact interval: the two base-four integer
parts differ by at most one.  Retain a deterministic subsequential
limit $r$ of this ratio and absorb it into the comparison constant.
Thus the $B(n)$ convention used in Section~\ref{comp:section} has
factor $r\lambda$, which we again denote by $\lambda$, with all
passage and circuit conclusions scaling together.  Neither an exact
ratio limit nor an additional time-scaling factor is required.
Uniform costs to prescribed
lattice endpoints are supplied separately in
Section~\ref{comp:section}.
\end{proof}
